\section{The best positive bound from prescribed samples}

Fix distinct offsets $a_1,\ldots,a_k>1/2$ and a target ordinate $t$. Write
$H_i=a_i\Re (\Lambda'_\chi/\Lambda_\chi)(1/2+a_i+it)$.
The preceding Hadamard calculation expresses $H_i$ as the sum over
$u=\gamma-t$ of $a_i^2/(a_i^2+u^2)$. Each admissible function
\[
 F(u)=\sum_{i=1}^k c_i\frac{a_i^2}{a_i^2+u^2},\qquad
 c_i\in\mathbb R,\quad F(0)=1,\quad F(u)\ge0\quad(u\in\mathbb R)
\]
gives the upper bound $m_t\le\sum c_iH_i$. The positivity condition
allows coefficients of either sign. Both the optimal certificate for a
given character and its asymptotic limitation can be computed exactly.

\paragraph{The optimal certificate.}
Form the real symmetric matrix
\[
 G_{ij}=\frac{a_jH_i+a_iH_j}{a_i+a_j},\qquad
 \mathbf1=(1,\ldots,1)^{\mathsf T}.
\]
Then $G$ is positive definite and the smallest admissible bound equals
\[
 \min_F\sum_\gamma F(\gamma-t)
    =\frac{1}{\mathbf1^{\mathsf T}G^{-1}\mathbf1}.
\]
Thus $G-2\mathbf1\mathbf1^{\mathsf T}$ failing to be positive
semidefinite is sufficient for quadratic central nonvanishing when $t=0$.

\emph{Proof.}
Use $g_i(u)=a_i/(a_i-iu)$, so $g_i(0)=1$, and symmetrize the zero
counting measure by giving mass $1/2$ at both $u$ and $-u$ for each
$u=\gamma-t$. This preserves the integral of any even function. Against
this measure the Gram matrix of the $g_i$ is real. The elementary identity
\[
 \Re\bigl(g_i(u)\overline{g_j(u)}\bigr)
 =\frac{a_j}{a_i+a_j}\frac{a_i^2}{a_i^2+u^2}
  +\frac{a_i}{a_i+a_j}\frac{a_j^2}{a_j^2+u^2}
\]
gives its entries. The integrals converge and positive definiteness
follows because a nonzero linear combination has finitely many zeros,
the poles are distinct, and the measure has infinite support.

Every admissible $F$ can be written $|\sum v_i g_i|^2$.
Indeed after multiplying by $\prod (u^2+a_i^2)$ its numerator is a real
polynomial of degree at most $2k-2$, nonnegative on $\mathbb R$.
Pair its conjugate nonreal roots and halve its real root multiplicities
to write it as $|P(u)|^2$ with $\deg P\le k-1$. Partial fractions
of $P(u)/\prod(a_i-iu)$ give the representation.
Multiplication by a phase makes $\sum v_i=1$.
Cauchy--Schwarz in the $G$ inner product gives the lower bound
$(\mathbf1^{\mathsf T}G^{-1}\mathbf1)^{-1}$, attained by
$v=G^{-1}\mathbf1/(\mathbf1^{\mathsf T}G^{-1}\mathbf1)$.
This vector is real, so its squared modulus is even. Its denominator
has distinct factors $u^2+a_i^2$ and its numerator has smaller degree,
hence the resulting function belongs to the stated span.
Finally a real vector with $v^{\mathsf T}(G-2\mathbf1\mathbf1^{\mathsf T})v<0$
has $\mathbf1^{\mathsf T}v\ne0$. Normalization then gives a bound below~2.

\paragraph{A sharp obstruction even with conductor-dependent coefficients.}
At $t=0$ set $S=\sum_i a_i^{-1}$ and hold the sites fixed. Uniformly over
primitive nonprincipal characters of conductor $q$, the minimum equals
\[
       \frac{\log q}{2S}+O_{a_1,\ldots,a_k}(1) \qquad(q\longrightarrow\infty).
\]
To see this write $G=(\log q)C+O(1)$ entrywise with
$C_{ij}=a_i a_j/(a_i+a_j)$. This uses absolute convergence of the
prime sums at all the sample sites. The matrix $C$ is positive definite since it is the Gram matrix of $x\mapsto a_i e^{-a_i x}$ on $[0,\infty)$.
For $w=C^{-1}\mathbf1$ put $R(z)=\sum_jw_j a_j/(z+a_j)$.
The equations $a_iR(a_i)=1$ give the rational identity
\[
 1-zR(z)=\prod_j\frac{a_j-z}{a_j+z}.
\]
Both sides have a common denominator, numerator degree at most $k$, the
same $k$ distinct zeros, and value $1$ at zero. Differentiation at zero
gives $\mathbf1^{\mathsf T}C^{-1}\mathbf1=2S$. Inversion of the
fixed-size matrix now yields the asymptotic. This rules out any
fixed choice of sampling sites as a uniform certificate, even if the
coefficients are selected optimally for each individual character.

The function attaining the least coefficient of $\log q$ is given by
$|R(-iu)/(2S)|^2$, or explicitly by
\[
 F_*(u)=\frac{\sin^2(\sum_j\arctan(u/a_j))}{S^2u^2},
 \qquad F_*(0)=1.
\]
For $k=1$ it is the single kernel $a_1^2/(a_1^2+u^2)$. For $k=2$ it is
the product of the two single kernels, as follows at once by the
sine addition formula. For $a_j=\sqrt j$ the optimal coefficient in
front of $\log q$ is $1/(2\sum_{j\le k}j^{-1/2})\sim 1/(4\sqrt{k})$.
The binomial test gives $c_k/2\sim1/(2\sqrt{\pi\log k})$.
These compare the leading coefficients for fixed sites as $q$ grows. The error bound depends on $k$.
The extremal calculation is an instance of a rational reproducing
kernel. For related classical theory see M.~M.~Dzhrbashyan,
\href{https://www.mathnet.ru/eng/sm3761}{\emph{Biorthogonal systems of
rational functions and best approximation of the Cauchy kernel on the
real axis}}, Math. USSR-Sb. 24 (1974), 409--433.

\paragraph{Allowing the sample sites to depend on the character.}
Let $M_k(\chi)$ be the infimum of the central bound over all distinct
$a_j>1/2$ and all admissible $F$ at those sites. For each fixed $k$,
uniformly over primitive nonprincipal characters as $q\to\infty$,
\[
                    \frac{M_k(\chi)}{\log q}\longrightarrow\frac{1}{4k}.
\]
Thus the obstruction persists even if both the sample locations and
the signed coefficients are selected after inspecting the character.

Here are details addressing sites that collide, escape to infinity, or
approach the edge of absolute convergence. Any such $F$ has the form
$P(u)/\prod_j(1+u^2/a_j^2)$, with $P$ even and nonnegative,
$P(0)=1$, and $\deg P\le 2k-2$. For fixed $b>1/2$ we consequently have
\[
 F(u)\ge(2b)^{-2k}f(u),\qquad
 f(u)=\frac{P(u)}{(1+u^2/b^2)^k},
\]
since each denominator factor is at most $1+4u^2\le
4b^2(1+u^2/b^2)$. The fixed-pole extremum above gives
$\int_{\mathbb R} f\ge\pi b/k$ also with this repeated pole. Indeed,
replace the $k$ occurrences of $b$ by distinct sites tending to $b$,
keep $P$ fixed, and use dominated convergence in the inequality
$\int F\ge\pi/(\sum_j1/a_j)$. The latter equality for the least
integral follows from
$\int a_j^2/(a_j^2+u^2)\,du=\pi a_j$
and the minimizing $C$-quadratic form above.

On this one fixed finite-dimensional space,
\[
 \left|\sum_\gamma f(\gamma)-
        \frac{\log q}{2\pi}\int_{\mathbb R} f(u)\,du\right|
      \le A_{b,k}\int_{\mathbb R}|f(u)|\,du .
\]
For the basis $(b^2+u^2)^{-j}$, $1\le j\le k$, this follows by
differentiating at $v=b^2$ the identity
\[
 \sum_\gamma\frac{1}{v+\gamma^2}
   =\frac{1}{\sqrt v}\Re\frac{\Lambda_\chi'}{\Lambda_\chi}
          (1/2+\sqrt v)
   =\frac{\log q}{2\sqrt v}+O_b(1).
\]
All derivatives through order $k-1$ of the error are uniformly bounded
in a real neighborhood of $b^2$. The differentiated prime-power series
converge absolutely there, and parity gives two possible gamma terms. The zero sums differentiate termwise
by local domination. Differentiating $\int_{\mathbb R}(v+u^2)^{-1}du
=\pi/\sqrt v$ yields the corresponding real integrals. Equivalence of
norms on the fixed finite-dimensional space proves the stated estimate.

Since $f\ge0$, taking the infimum of the lower bound and then $q\to\infty$
gives $\liminf M_k(\chi)/\log q\ge(2b)^{-2k}b/(2k)$.
Letting $b\downarrow1/2$ gives $1/(4k)$. Conversely take fixed distinct
$a_j$ arbitrarily close to $1/2$ from above and apply the fixed-site
minimum, whose leading coefficient tends to $1/(4k)$. This proves the
matching upper bound. The limit is for $k$ fixed and makes no assertion
about the regime $k=k(q)\to\infty$.
